% !TEX root = ../../root.tex
A recent state-space framework based on Partial Integral Equations (PIEs)
provides a unified formulation of stabilization and
$H_\infty$-optimal state-feedback controller synthesis for PDEs
\cite{shivaknew}. This framework addresses the nominal setting, corresponding
to $\Delta=0$. Building on this result, we develop a unified framework for
optimal state-feedback control of uncertain or nonlinear PIEs subject to a
prescribed quadratic performance requirement. As established in the previous
section, IQCs are used to characterize the uncertainty and reduce conservatism
in the resulting robust-stability certificates. Let $G$ denote the open-loop
PIE plant to which the controller is connected. Its realization is given in
Eq.~\eqref{eq:primal-system-components}. For a controller $K$, define the
closed-controller generalized plant as $P_K:=K\star G$.

Our objective is to construct a state-feedback controller $K$ such that
$P_K\star\Delta$ is robustly stable and satisfies the prescribed performance
IQC for every $\Delta\in\mbf\Delta$. We consider the controller
interconnection shown in Fig.~\ref{fig:primal-controller-interconnection}, with
\begin{equation}\label{eq:primal-controller}
    u
    = \mcl K_{\mrm G}\mbf x_{\mrm G}
    + \mcl K_{\Theta}\mbf x_{\Theta}.
\end{equation}
The auxiliary filter $\Theta$ and its associated state $\mbf x_{\Theta}$ are
constructed below. For dynamic IQCs, causal implementation of the
state-feedback law requires that the signals $w$ and $z$ driving the dynamic
part of $\Theta$ be available to the controller, so that $\mbf x_{\Theta}$ can
be computed in real time.

\primalcontrollerfigure

\subsection{Inputs at the boundary}

We also admit boundary control and disturbance inputs through the treatment
introduced by Shivakumar et al.~\cite[Sec.~III-B]{shivaknew}. The following
construction and heat-equation example are adapted directly from that work.
Let $x(t,s)=\partial_s^2v(t,s)$ and impose the boundary conditions
$v(t,0)=u_1(t)$ and $\partial_sv(t,0)=u_2(t)$. The distributed state is
reconstructed as
\begin{equation}\label{eq:boundary-input-reconstruction}
v(t,s)=u_1(t)+su_2(t)
+\int_0^s(s-\eta)x(t,\eta)\,\mrm d\eta.
\end{equation}
For the heat equation $\partial_tv=\partial_s^2v$, this gives
\begin{equation}\label{eq:boundary-input-pie}
\partial_t\!\int_0^s(s-\eta)x(t,\eta)\,\mrm d\eta
=x(t,s)-\dot u_1(t)-s\dot u_2(t).
\end{equation}
Thus a boundary signal enters the PIE through its time derivative. Retaining
this derivative while measuring the gain from the undifferentiated signal in
$L_2$ would require additional temporal regularity. Following
\cite{shivaknew}, for a generic boundary signal $w_{\mrm b}$ we therefore set
$\hat w_{\mrm b}=\dot w_{\mrm b}$ and regard $\hat w_{\mrm b}$ as the PIE
input.

If $w_{\mrm b}$ is required by the controller or an IQC filter, it is
recovered through
\begin{equation}\label{eq:boundary-input-integrator}
\dot w_{\mrm b}(t)=\hat w_{\mrm b}(t),
\qquad
w_{\mrm b}(t)=w_{\mrm b}(0)
+\int_0^t\hat w_{\mrm b}(\tau)\,\mrm d\tau.
\end{equation}
We assume that $w_{\mrm b}(0)$ is known and that $w_{\mrm b}$ is available
wherever it is used. For the zero-state performance problem,
$w_{\mrm b}(0)=0$, consistently with the initial boundary value of the PDE.
The integrator state is included in the augmented plant realization. Hence the
system class includes distributed inputs and boundary control or disturbance
inputs for which this augmentation produces a well-posed PIE realization.
Direct uncertainty through a boundary condition is outside this class unless
it can also be represented by a well-posed PIE uncertainty channel satisfying
the assumed hard IQC.

\subsection{Primal analysis and synthesis obstruction}

To distinguish analysis from synthesis, fix the IQC filter $\Psi$ and write an
augmented realization before closing the controller as
\begin{equation}\label{eq:primal-filtered-open-loop}
\bmat{\mcl T\dot{\mbf x}\\\tilde v}
=
\bmat{\mcl A&\mcl B_{\mrm u}&\mcl B\\
       \mcl C&\mcl D_{\mrm{zu}}&\mcl D}
\bmat{\mbf x\\u\\w}.
\end{equation}
With $u=\mcl K\mbf x$, where
$\mcl K=\bmat{\mcl K_{\mrm P}&\mcl K_\Theta}$, the filtered graph
$\Psi\bmat{K\star G\\I}$ has the PIE realization
\[
\{\mcl T,\mcl A+\mcl B_{\mrm u}\mcl K,\mcl B,
\mcl C+\mcl D_{\mrm{zu}}\mcl K,\mcl D\}.
\]
A sufficient condition for robust stability and performance is the existence
of a PI operator $\mcl P=\mcl P^*\succeq0$ and $\epsilon>0$ satisfying
\begin{equation}\label{eq:primal-controller-analysis}
\begin{gathered}
\bmat{\mcl T^*\mcl P(\mcl A+\mcl B_{\mrm u}\mcl K)
+(\mcl A+\mcl B_{\mrm u}\mcl K)^*\mcl P\mcl T
&\mcl T^*\mcl P\mcl B\\
\mcl B^*\mcl P\mcl T&\epsilon I}\\
+\bmat{(\mcl C+\mcl D_{\mrm{zu}}\mcl K)^*\\\mcl D^*}
\mcl V
\bmat{(\mcl C+\mcl D_{\mrm{zu}}\mcl K)&\mcl D}
\preceq0.
\end{gathered}
\end{equation}
For analysis, $\mcl K$ is fixed and
\eqref{eq:primal-controller-analysis} is affine in the certificate operators
$\mcl P$ and $\mcl V$. For synthesis, these operators and $\mcl K$ must be
determined simultaneously. This introduces products between $\mcl P$ and
$\mcl K$, and between $\mcl K$ and $\mcl V$.
For a non-descriptor ODE, the storage and controller coupling can be removed
with $\mcl Q=\mcl P^{-1}$ and $\mcl Z=\mcl K\mcl Q$. The reconstruction
operator $\mcl T$ prevents the same cancellation for a general PIE. The
multiplier coupling must likewise be resolved by additional structure. Thus
the primal formulation provides a convex test for a fixed controller, but not
a convex synthesis condition. The following sections construct a dual
equivalent of the filtered interconnection from which $K$ can be synthesized
convexly.
